\clearpage

\section {Analysis procedure}
\label{sec:oscfit}

%--------------------------------------------------------------------------------------------------------------------
\subsection {Fit method}
\label{sec:oscfit_method}

The objective of this analysis is to search for $\bar{\nu}_e$ appearance from a $\bar{\nu}_{\mu}$ beam.
The parameter of interest is $\beta$, a not physical parameter which was introduced to normalize the $\bar{\nu}_e$ appearance signal component
(please see figure~\ref{fig:1Re_spectrum_and_components_default_osc_postfit}).
The hypothesis of null $\bar{\nu}_e$ appearance is obtained by setting $\beta = 0$, 
while $\bar{\nu}_e$ appearance corresponding to the current best knowledge of the PMNS matrix is given by $\beta = 1$.
All the systematic parameters are marginalized.
Only the solar parameters 
$\sin^2 \theta_{12}$ 
and 
$\Delta m^2_{21}$ 
are fixed because they have a negligible 
impact on the oscillations at the T2K energy. 
The other oscillation parameters are marginalized using the prior distributions which correspond to 
the posterior of the T2K+reactor joint fit bayesian analysis, which are shown in section~\ref{sec:systematics:inputs}. 
The treatment of the oscillation parameters is shown in table~\ref{tab:osc_params_fixed_values_reactor}.

%--------------------------------------------------------------------------------------------------------------------

The {\em VALOR} appearance analysis uses a binned likelihood-ratio method where
the estimation of the $\beta$ parameter is obtained
by comparing the observed and predicted SuperK reconstructed energy spectra
for 1-ring $e$-like events.

Let $N$ be the number of reconstructed energy bins and
${\bf a}$ be a $1 \times N_{s}$ - dimensional array of systematic (nuisance) 
parameters that affect the SuperK reconstructed energy spectrum prediction (e.g.
parameters that control the beam, neutrino interaction simulations, 
detector response and oscillations).
The log-likelihood function used in the analysis is
\begin{equation}
\displaystyle
%\chi^{2} 
-2 \; \textrm{ln} \; \lambda (\beta,{\bf a}) = 2 \cdot \sum_{i=0}^{N-1} 
   \Big( 
     n^{obs}_{i} \cdot \textrm{ln} (n^{obs}_{i}/n^{exp}_{i}) +
     (n^{exp}_{i}-n^{obs}_{i}) 
   \Big) +
   ({\bf a}-{\bf a_{0}})^{T} \cdot {\bf C}^{-1} \cdot ({\bf a}-{\bf a_{0}})
\label{eq:chisquare}
\end{equation}
where
${\bf a_{0}}$ is an $1 \times N_{s}$ - dimensional array with the default values of 
the systematic parameters, ${\bf a}^{T}$ is the transpose of ${\bf a}$,
${\bf C}$ is the systematic parameter covariance matrix of dimension $N_{s} \times N_{s}$,
$n^{obs}_{i}$ is the observed number of events in the $i^{th}$ bin, and
$n^{exp}_{i} = n^{exp}_{i}(\beta; {\bf a})$ 
%$n^{exp}_{i} = n^{exp}_{i}(\sin^{2}\overline{\theta}_{23},|\overline{\Delta m}^{2}_{32}|; {\bf a})$ 
is the corresponding expected number of events.

The advantage of the likelihood ratio method over the extended maximum-likelihood method is that,
in the large-sample limit, the quantity 
$-2 \: \textrm{ln} \: \lambda (\beta; {\bf a})$
%$-2 \: \textrm{ln} \: \lambda (\sin^{2}\overline{\theta}_{23},|\overline{\Delta m}^{2}_{32}|; {\bf a})$
in Eq.\ref{eq:chisquare} has a $\chi^{2}$ distribution and it can therefore be used as a goodness-of-fit test.

Values of systematic parameters that give a negative predicted number of events in any reconstructed energy bin in any interaction 
mode are not allowed. If a systematic throw gives a negative number of events, that toy experiment is rejected.

%--------------------------------------------------------------------------------------------------------------------                                                                                    

\subsection {Marginalization}
\label{sec:marginalization}

In the previous analyses a fully frequentist approach, where all the nuisance parameters are profiled, was used.
Now a hybrid bayesian-frequentist method is used. The nuisance parameters are marginalized while the parameter of interest is minimized.
The marginalization of the nuisances gives a so-called marginal likelihood that depends only on $\beta$:

\begin{equation}
\lambda_{marg} (\beta) = \int_{A} \lambda (\beta;{\bf a}) \pi{({\bf a})} d {\bf a} = \sum_{i=0}^{n} \lambda (\beta; {\bf a}_i) \pi({ {\bf a}_i})  
\label{eq:likemarg}
\end{equation}

where $n$ is the number of toys, ${\bf a}$ is a single extracted set of nuisance parameter, $\pi{({\bf a})}$ is the prior of the nuisance parameters and 
$\pi( {\bf a}_i)$ is the weight of a particular set of nuisance parameters (e.g. $w=1/n$ if flat prior).
Many sets of ${\bf a}$ ($n = 10k$), should be extracted in order to sample the space of the nuisance parameters.
The number of extractions depends on the number of nuisance parameters as well as the precision that is required by the measurement.

We derive the marginal log-likelihood, which we will denoted as $\chi^2$ from now on:

\begin{equation}
%\chi^2 (\beta) = - 2 \textrm{\ln} \lambda_{marg} (\beta)  = -2 \textrm{\ln} \left( \sum_{i=0}^{n} \lambda ({\bf \beta}; {\bf a}_i) \pi({\bf a_i})  \right) 
\chi^2 (\beta) = - 2  \ln \lambda_{marg} (\beta)  = -2 \ln \left( \sum_{i=0}^{n} \lambda ({\bf \beta}; {\bf a}_i) \pi({\bf a_i})  \right) 
\label{eq:loglikemarg}
\end{equation}

It is important to note that this numerical integration is a weighted average of the likelihood evaluated at the $n$ throws, 
not a weighted average of the log-likelihood.
When marginalization is used the penalty term
$ ({\bf a}-{\bf a_{0}})^{T} \cdot {\bf C}^{-1} \cdot ({\bf a}-{\bf a_{0}}) $
becomes constant, independent of the parameter of interest.
For this reason a modified version of eq.~\ref{eq:chisquare}, which doesn't include the penalty term, is used.

The toys used for the marginalization of the nuisance parameters are performed exactly in the same way as 
described in 
sec.~\ref{sec:systematics:inputs}
, taking into account the correlations, but without statistical fluctuations.
First $n$ reconstructed energy distributions ($E_{reco}^{marg}$) are produced using the randomized systematic parameters
and normalized to $1 \times 10^{20}$ POT, in order to be used for fits at any different number of POT. 
Finally they are stored in output files which can be used for whatever needed fit.
A fit on a single data set is performed as following:
\begin{itemize}
\item all the $n$ toy reconstructed energy distributions ($E_{reco}^{marg}$) are read from the files,
\item each $E_{reco}^{marg}$ distribution is normalized to the number of POT of the data set,
\item for each $E_{reco}^{marg}$ eq.~\ref{eq:chisquare} is calculated and $\lambda(\beta,{\bf a}) = e^{- (-2 \; \textrm{ln} \; \lambda (\beta,{\bf a})) / 2}$ is obtained,
\item the $\chi^2$ value is finally obtained, using eq.~\ref{eq:loglikemarg}, for a given value of $\beta$. 
\end{itemize}


%--------------------------------------------------------------------------------------------------------------------                                                                                    

\subsection {P-value}
\label{sec:pvalue}

% Describe the p-values used for the analysis

There are different test statistics which can be used to quote a p-value,
as described in 
~\cite{slides_davide_pvalue_asg} % ASG talk
.
As proposed by the OA group, for this $\bar{\nu}_e$ appearance search two different p-values are quoted.
In order to quote the significance given null $\bar{\nu}_e$ appearance, the parameter of interest is set to $\beta=0$
and $n=10k$ toys are performed.
The systematic parameters are randomized taking into account the correlations as described in sec.~\ref{sec:systematics:inputs}.
Then the reconstructed energy distribution is fluctuated with Poisson. 
Then the following p-values are calculated:
\begin{itemize}

\item p-value with rate-only analysis: the number of events is used as test statistics, without using the shape information of the reconstructed energy 
	distribution. The distribution $f(N_{SK})$ of the total number of superK events in a toy ($N_{SK}$) is obtained from the performed toys. 
	Then the number of events is measured from the data set ($N_{SK}^{obs}$) and
	the 1-sided (right tail) p-value is computed in the following way: 	
	
	\begin{equation}
	\text{p-value (right tail)} = \int_{N_{SK}^{obs}}^{+\infty} f(N_{SK}) d N_{SK}  = 
		\frac{ \sum_{ itoy=1 }^{n} w(N_{SK} \geq N_{SK}^{obs} )^{itoy} }{ n }
	\label{eq:pvalue_rateonly}
	\end{equation}
	
	where $w(N_{SK} \geq N_{SK}^{obs})$ is 1 if $N_{SK} \geq N_{SK}^{obs}$ otherwise 0.
	The distribution $f(N_{tot}^{toy})$ is expected to be discrete since it follows the Poisson distribution.
			 
	In an analogous way the left tail 1-sided p-value is computed as following:
		
	\begin{equation}
	\text{p-value (left tail)} = \int^{N_{SK}^{obs}}_{-\infty} f(N_{SK}) d N_{SK}  = 
		\frac{ \sum_{ itoy=1 }^{n} w(N_{SK} \leq N_{SK}^{obs} )^{itoy} }{ n }
	\label{eq:pvalue_rateonly}
	\end{equation}
		 
\item p-value with rate+shape analysis: the shape information of the reconstructed energy distribution is used.
	In this case the test statistic is used to compare the case of $\beta=0$ with $\beta=1$:
	
	\begin{equation}
	\Delta \chi^2 = \chi^2(\beta=0) - \chi^2(\beta=1)
	\label{eq:deltachisq_rateshape}
	\end{equation}

	where $\chi^2(\beta=0)$ and $\chi^2(\beta=1)$ are computed using eq.~\ref{eq:loglikemarg}.
	It is worth to note that the parameter $\beta$ is not fitted but fixed either to 0 or 1.
	The 1-sided p-value (right tail) is calculated as:

	\begin{equation}
	\text{p-value (right tail)} = \int_{\Delta \chi^2_{obs}}^{+\infty} f(\Delta \chi^2) d \Delta \chi^2  = 
		\frac{ \sum_{ itoy=1 }^{n} w(\Delta \chi^2 \geq \Delta \chi^2_{obs} )^{itoy} }{ n }
	\label{eq:pvalue_rateonly}
	\end{equation}
	
	where $w(\Delta \chi^2  \geq \Delta \chi^2 _{obs})$ is 1 if $\Delta \chi^2  \geq \Delta \chi^2 _{obs}$ otherwise 0.

	The 1-sided p-value (left tail) is calculated as:

	\begin{equation}
	\text{p-value (left tail)} = \int^{\Delta \chi^2_{obs}}_{-\infty} f(\Delta \chi^2) d \Delta \chi^2  = 
		\frac{ \sum_{ itoy=1 }^{n} w(\Delta \chi^2 \leq \Delta \chi^2_{obs} )^{itoy} }{ n }
	\label{eq:pvalue_rateonly}
	\end{equation}
	
\end{itemize}

The p-value will be calculated on the left or right tail depending on the position of the observed test statistic
in the expected distribution. For example, if the observed test statistic is on the left tail of the expected distribution of the
test statistic, a 1-sided p-value is computed on the left tail. Otherwise a 1-sided p-value on the right tail is computed. 

%--------------------------------------------------------------------------------------------------------------------                                                                                    


